\begin{tabular}{rlrrrp{0.5\linewidth}}
\toprule
\textbf{\textbackslash\{\}\#} & \textbf{True} & \textbf{p(AI)} & \textbf{Size} & \textbf{Sharpness} & \textbf{Hypothesis (from measured cues)} \\
\midrule
1 & AI & 0.05 & 1200x800 & 2699 & very sharp, heavy detail or noise looks 'camera-like' \\
2 & AI & 0.06 & 6000x4001 & 4353 & very sharp, heavy detail or noise looks 'camera-like'; highly saturated, painted / illustration style \\
3 & Real & 0.92 & 1920x1080 & 2115 & very sharp, heavy detail or noise looks 'camera-like'; extreme aspect ratio squashed by resizing \\
4 & AI & 0.08 & 1024x512 & 609 & extreme aspect ratio squashed by resizing \\
5 & AI & 0.11 & 1200x600 & 2290 & very sharp, heavy detail or noise looks 'camera-like'; extreme aspect ratio squashed by resizing \\
6 & Real & 0.86 & 681x960 & 1517 & very sharp, heavy detail or noise looks 'camera-like'; highly saturated, painted / illustration style \\
7 & Real & 0.85 & 2880x1620 & 2290 & very sharp, heavy detail or noise looks 'camera-like'; extreme aspect ratio squashed by resizing \\
8 & Real & 0.84 & 390x280 & 1448 & low resolution: up-scaling to the input size hides texture cues; very sharp, heavy detail or noise looks 'camera-like' \\
9 & Real & 0.80 & 1200x690 & 3407 & very sharp, heavy detail or noise looks 'camera-like'; grey-scale image: no colour statistics to rely on \\
10 & AI & 0.21 & 1200x800 & 3447 & very sharp, heavy detail or noise looks 'camera-like' \\
\bottomrule
\end{tabular}
